\documentclass[11pt]{article}

\usepackage{amsmath,amssymb,bm,booktabs,graphicx,hyperref}

\title{Trial NNN: Short Descriptive Title}
\author{}
\date{}

\begin{document}
\maketitle

% Copy this file to gp/results/trial_NNN/report.tex and fill in each section.
% Compile with: pdflatex report.tex

\section{Governing Physics}

% Level-set evolution equation aligned with simulation.ipynb.
\begin{equation}
  \label{eq:levelset}
  % \phi_t + R_f(\phi)\,|\nabla\phi| - \nu\,\Delta\phi = 0
\end{equation}

% Rothermel rate-of-spread terms embedded in the RHS.

\begin{table}[h]
  \centering
  \caption{Constants and parameters (from \texttt{simulation.ipynb}).}
  \begin{tabular}{@{}lll@{}}
    \toprule
    Symbol & Value & Description \\
    \midrule
    % $\Delta x$ & & grid spacing \\
    \bottomrule
  \end{tabular}
\end{table}

\section{Method}

\subsection{Prior placement}

% Where the GP prior sits and why.
\begin{equation}
  \label{eq:prior}
  % \phi^{n-1}(\mathbf{x}) \sim \mathcal{GP}\bigl(\mu(\mathbf{x}),\, k(\mathbf{x},\mathbf{x}';\theta)\bigr)
\end{equation}

\subsection{Euler discretization and numerical GP construction}

% Forward Euler update.
\begin{equation}
  \label{eq:euler}
  % \phi^n = \phi^{n-1} + \Delta t\,\mathcal{L}_x \phi^{n-1}
\end{equation}

% Joint multi-output GP structure (block matrix notation).
\begin{equation}
  \label{eq:joint-gp}
  \begin{bmatrix} \phi^n \\ \phi^{n-1} \end{bmatrix}
  \sim \mathcal{GP}\left(
    0,\,
    \begin{bmatrix}
      k^{n,n} & k^{n,n-1} \\
      k^{n-1,n} & k^{n-1,n-1}
    \end{bmatrix}
  \right)
\end{equation}

% Kernel expressions derived from the time-stepper.
% If only the posterior mean is propagated, state the equations and what covariance terms are dropped.

\subsection{Inference from sparse observations}

% Observation model, 50-point sampling setup (seed, locations).
\begin{equation}
  \label{eq:posterior}
  % \mu_* = \mu + K_{*X}(K_{XX} + \sigma_n^2 I)^{-1}(y - \mu_X)
\end{equation}

% Kernel choice, hyperparameters, and fitting procedure.

\subsection{Algorithm}

\begin{enumerate}
  \item Sample sparse observations at $t=0$.
  \item Infer latent initial field via GP conditioning.
  \item Advance with Euler numerical GP update to each evaluation time.
  \item Compare prediction to \texttt{simulation.ipynb} ground truth.
\end{enumerate}

\section{Assumptions and approximations}

\begin{itemize}
  \item % e.g.\ mean-only propagation; ignored boundaries; nonlinear term evaluated at posterior mean
\end{itemize}

\section{Experimental setup}

% Seed, grid size, $\Delta t$, $\Delta x$, evaluation times, code paths.

\section{Results}

\begin{table}[h]
  \centering
  \caption{Metrics across evaluation times.}
  \begin{tabular}{@{}rcccc@{}}
    \toprule
    Time & Field RMSE & Field MAE & Front distance & Notes \\
    \midrule
    % 50 & & & & \\
    \bottomrule
  \end{tabular}
\end{table}

% Reference figure filenames, e.g.:
% \begin{figure}[h]
%   \centering
%   \includegraphics[width=0.9\linewidth]{comparison_t100.png}
%   \caption{Predicted vs.\ ground truth at $t=100$.}
% \end{figure}

\section{Sources}

\begin{itemize}
  \item Raissi et al., \emph{Numerical Gaussian Processes for Time-dependent and Non-linear PDEs}.
  \item Rothermel (1972); WRF level-set fire-spread paper (\texttt{gp/docs/}).
  \item \texttt{simulation.ipynb} (implementation reference).
\end{itemize}

\section{Notes (optional)}

\begin{itemize}
  \item % what worked, what failed, ideas for next trial
\end{itemize}

\end{document}
